\documentclass[11pt]{paper}

\usepackage[T1]{fontenc}

\usepackage{graphicx}
\usepackage{amssymb}
\usepackage{amstext}
\usepackage{amsmath}
\usepackage{a4wide,color}
\usepackage[utf8]{inputenc}
\usepackage[frenchb]{babel}
\usepackage{xspace}
\usepackage{anysize}
\usepackage{tabularx}
\usepackage{multirow}
\usepackage{fancybox}
\usepackage{fancyhdr}
\usepackage{bbding}

\usepackage{siunitx}

\usepackage{color}
\usepackage{float}

\usepackage[toc,page]{appendix} 
\usepackage{lscape}

\usepackage{placeins}

\usepackage{listings}
\lstdefinestyle{customc}{
  belowcaptionskip=1\baselineskip,
  breaklines=true,
  frame=L,
  xleftmargin=\parindent,
  language=C,
  showstringspaces=false,
  basicstyle=\tiny\ttfamily,
  keywordstyle=\bfseries\color{green!40!black},
  commentstyle=\itshape\color{purple!40!black},
  identifierstyle=\color{blue},
  stringstyle=\color{orange},
  tabsize=2
}

\usepackage{todonotes}

\newcommand{\version}{1.0.1{ }}

\usepackage[framemethod=TikZ]{mdframed}

\lstset{% general command to set parameter(s)
	basicstyle=\footnotesize, % print whole listing small
	keywordstyle=\color{magenta}\bfseries, % underlined bold black keywords
	identifierstyle=, % nothing happens
	commentstyle=\color{black}, % white comments
	stringstyle=\ttfamily, % typewriter type for strings
	showstringspaces=false % no special string spaces
} 

\lstdefinelanguage{aald}
	{morekeywords={system, implementation},
	sensitive=false,
	morecomment=[l]{//},
	morecomment=[s]{/*}{*/},
	morestring=[b]",
}

\usepackage[colorlinks=true]{hyperref}

\usepackage[normalem]{ulem}
\usepackage{color}

\definecolor{Fond}{gray}{0.7}

\renewcommand{\floatpagefraction}{.99}
\renewcommand{\textfraction}{.01}
\newcommand{\modif}[1]{\textcolor{red}{\uline{#1}}}


\pagestyle{fancy}
\fancyhf{}
\fancyhead[RE,RO]{\thepage}
\fancyhead[LE]{}
\fancyhead[LO]{Temps réel }
\fancyfoot[LO]{Centrale Nantes -- \today}
\fancypagestyle{plain}{%
  \fancyhf{} % get rid of headers
  \renewcommand{\headrulewidth}{0pt} % and the line
}

\newenvironment{maliste}%
{ \begin{list}%
	{\ArrowBoldRightStrobe}%
	{\setlength{\labelwidth}{30pt}%
	 \setlength{\leftmargin}{35pt}%
	 \setlength{\itemsep}{\parsep}}}%
{ \end{list} }

 \setlength{\parskip}{1.2ex plus .9ex minus .4ex}

\newcommand{\ocv}{{\tt OpenCV}\xspace} 
\newcommand{\alsa}{{\tt ALSA}\xspace} 
\newcommand{\aadl}{{\tt AADL}\xspace} 
\newcommand{\xeno}{{\tt Xenomai}\xspace} 
\newcommand{\iic}{I$^2$C\xspace} 


\title{{\Large Systèmes interconnectés -- SINT }\\
\vspace{-0.4em}
{\large Sujet de TP 3B : BONUS}\\
{\small Centrale Nantes}\\
\small P.-E. Hladik, \small{\texttt{pierre-emmanuel.hladik@ec-nantes.fr}}\\
---\\
{\scriptsize  Version \version (\today)}\\
}


\begin{document}

\maketitle

%\begin{table}[htdp]
%\caption{Suivi des modifications}
%\begin{center}
%\begin{tabular}{|c|c|c|p{8cm}|}
%\hline
%Date & Version & Auteur & Modifications\\
%\hline
%JJ/MM/12 & 0.3.X & P.-E. Hladik&  \\
%\hline
%\end{tabular}
%\end{center}
%\label{default}
%\end{table}%

%%%%%%%%%%%%%%%%%%%%%%%%%
  \section{Objectifs pédagogiques}
%%%%%%%%%%%%%%%%%%%%%%%%%


\begin{itemize}
	\item intégrer les différents protocole sur une même carte
\end{itemize}


%%%%%%%%%%%%%%%%%%%%%%%%%
\section{Cahier des charges}
%%%%%%%%%%%%%%%%%%%%%%%%%

\subsection{Objet}
Concevoir une application embarquée permettant de \textbf{configurer l'accéléromètre} à l'aide de \textit{boutons} et d'un \textit{encodeur rotatif avec bouton}, d'afficher l'état courant sur l'écran, puis de lancer un \textbf{mode acquisition} qui affiche et enregistre les mesures des trois axes dans la mémoire du \texttt{MCP2515}.

\subsection{Modes de fonctionnement}
\begin{itemize}
  \item \textbf{Mode normal} : l'application est en attente d'une action utilisateur.
  \item \textbf{Mode configuration} : sélection du plein-échelle de l'accéléromètre (\SI{2}{g}, \SI{4}{g}, \SI{8}{g}).
  \item \textbf{Mode acquisition} : lecture continue des axes X, Y, Z, affichage à l'écran et écriture en mémoire du \texttt{MCP2515}.
\end{itemize}

\subsection{Entrée en mode configuration}
\begin{itemize}
  \item \textbf{Déclenchement} : appui simultané sur \textbf{Bouton~1} et \textbf{Bouton~2}.
  \item \textbf{Effet immédiat} :
    \begin{itemize}
      \item Interrompre toute mesure/acquisition en cours.
      \item Afficher explicitement sur l'écran : \texttt{Mode : CONFIGURATION}.
    \end{itemize}
\end{itemize}

\subsection{Paramétrage du plein-échelle}
\begin{itemize}
  \item \textbf{Sélection} : l'encodeur rotatif permet de parcourir les options de plein-échelle suivantes : \SI{2}{g}, \SI{4}{g}, \SI{8}{g}.
  \item \textbf{Affichage} : l'écran doit refléter en temps réel l'option surlignée, p.~ex. \texttt{Plein-échelle : 4g (non validé)}.
  \item \textbf{Validation} : un appui sur le \textbf{bouton de l'encodeur} valide l'option actuellement sélectionnée.
  \item \textbf{Retour visuel} : afficher un message de confirmation, p.~ex. \texttt{Configuration validée : 4g}.
  \item \textbf{Application} : la configuration validée est immédiatement appliquée à l'accéléromètre.
\end{itemize}

\subsection{Lancement du mode acquisition}
\begin{itemize}
  \item \textbf{Déclenchement} : appui sur le \textbf{Bouton~3}.
  \item \textbf{Affichage} : afficher \texttt{Mode : ACQUISITION}.
  \item \textbf{Affichage des mesures} : montrer en continu les trois valeurs courantes :
        \begin{center}
          \texttt{X = \dots,\quad Y = \dots,\quad Z = \dots}
        \end{center}
  \item \textbf{Enregistrement} : écrire les trois valeurs (X, Y, Z) dans la mémoire du \texttt{MCP2515} à chaque acquisition.
\end{itemize}

\subsection{Exigences d'affichage (exemples indicatifs)}
\begin{itemize}
  \item \texttt{Mode : CONFIGURATION} lorsqu'on est dans l'écran de réglage.
  \item \texttt{Plein-échelle : 2g/4g/8g} pendant la navigation à l'encodeur.
  \item \texttt{Configuration validée : Xg} après validation par le bouton de l'encodeur.
  \item \texttt{Mode : ACQUISITION} lorsqu'on enregistre et affiche les mesures.
  \item \texttt{X = \dots\ Y = \dots\ Z = \dots} en acquisition.
\end{itemize}

\subsection{Critères d'acceptation}
\begin{enumerate}
  \item Un appui simultané sur \textbf{Bouton~1 + Bouton~2} interrompt l'acquisition (si en cours) et affiche \texttt{Mode : CONFIGURATION}.
  \item La rotation de l'encodeur fait défiler \SI{2}{g}, \SI{4}{g}, \SI{8}{g}, et l'écran reflète toujours la sélection courante.
  \item Un appui sur le \textbf{bouton de l'encodeur} confirme le plein-échelle et affiche \texttt{Configuration validée : Xg}.
  \item Un appui sur \textbf{Bouton~3} affiche \texttt{Mode : ACQUISITION}, montre en continu X, Y, Z et écrit ces trois valeurs dans la mémoire du \texttt{MCP2515}.
\end{enumerate}




\end{document}